\chapter{Prompts and Schema}
\label{app:prompts}

\section{JSON schema given to the model}
\begin{lstlisting}[language=JSON, caption={Schema embedded in every generation prompt}, label={lst:schema}]
{
  "mcq": [{"question": str, "options": [str, str, str, str],
           "answer_index": 0-3, "explanation": str, "source_quote": str}],
  "short": [{"question": str, "model_answer": str, "key_points": [str]}],
  "coding": [{"question": str, "language": str, "starter_code": str,
              "solution": str,
              "test_cases": [{"input": str, "expected_output": str}],
              "explanation": str}]
}
\end{lstlisting}

\section{Difficulty guidance}
\begin{table}[H]
\centering
\begin{tabularx}{\textwidth}{@{} p{2cm} X @{}}
\toprule
\textbf{Level} & \textbf{Guidance placed in the prompt} \\
\midrule
Easy & Recall and definitions; direct facts from the text; obvious distractors. Coding: simple one-concept tasks (loops, conditionals, basic functions). \\
Medium & Understanding and application; small scenarios; plausible distractors. Coding: moderate logic, standard data structures, small algorithms. \\
Hard & Analysis and problem solving; multi-step reasoning, edge cases, subtle distractors. Coding: optimisation, complexity analysis, tricky edge cases. \\
\bottomrule
\end{tabularx}
\caption{Difficulty guidance strings}
\label{tab:difficulty}
\end{table}

\section{System prompts}
\begin{lstlisting}[style=pseudo, caption={Document-quiz system prompt (abridged)}, label={lst:sys-doc}]
You are an expert university examiner for B.Tech Computer Science.
Create quiz questions ONLY from the provided material. For each MCQ,
ensure distractors are plausible and include a source_quote (exact
short excerpt justifying the answer). Respond with a single valid JSON
object and nothing else (no markdown fences).
\end{lstlisting}

\begin{lstlisting}[style=pseudo, caption={Topic-quiz system prompt (abridged)}, label={lst:sys-topic}]
You are an expert university examiner and subject tutor. Create
accurate, well-balanced quiz questions about the requested topic. When
web research material is provided, base questions on it and quote it in
source_quote; otherwise use well-established, uncontroversial knowledge.
Never invent facts. Respond with a single valid JSON object only.
\end{lstlisting}
